\documentclass[11pt,letterpaper]{article}
\usepackage{iftex}
\ifPDFTeX\usepackage[utf8]{inputenc}\fi
\usepackage[T1]{fontenc}
\usepackage[spanish,es-nodecimaldot]{babel}
\usepackage[margin=2cm,headheight=15pt]{geometry}
\usepackage{lmodern,microtype}
\usepackage{amsmath}
\usepackage[table]{xcolor}
\usepackage{array,booktabs,longtable,tabularx}
\usepackage{enumitem,fancyhdr,hyperref}
\usepackage{tikz}
\definecolor{marino}{HTML}{002D62}
\definecolor{dorado}{HTML}{F2A900}
\definecolor{suave}{HTML}{EDF3F8}
\hypersetup{colorlinks=true,linkcolor=marino,urlcolor=marino}
\pagestyle{fancy}
\fancyhf{}
\fancyhead[L]{\small EmprendeBúho}
\fancyhead[R]{\small Ingeniería de Software II}
\fancyfoot[C]{\thepage}
\setlength{\parindent}{0pt}
\setlength{\parskip}{6pt}
\setlength{\emergencystretch}{2em}
\setlength{\tabcolsep}{4pt}
\renewcommand{\arraystretch}{1.17}
\setlist[itemize]{nosep,leftmargin=*}
\newcolumntype{L}[1]{>{\raggedright\arraybackslash}p{#1}}
\newcolumntype{Y}{>{\raggedright\arraybackslash}X}
\newcommand{\hu}[1]{\mbox{HU-#1}}
\newcommand{\terminada}{\textbf{Terminada}}
\newcommand{\pendiente}{\textbf{Pendiente}}
\newcommand{\enproceso}{\textbf{En proceso}}
\newcommand{\notaBorrador}{\begingroup\small
\textbf{Nota de elaboración.} Reconstrucción propuesta a partir de las bitácoras
1 y 2, el primer avance, el cronograma original y el estado global comunicado
por el equipo. La distribución diaria de actividades y cierres requiere
validación de sus responsables. La semana 6 se presenta completa para revisión
al 13 de noviembre; las entradas posteriores al 9 de noviembre son previstas.
Los estados documentales no constituyen una auditoría del código ni un acta de
pruebas.\par\endgroup}
\newcommand{\portada}[2]{%
\hypersetup{pageanchor=false}\begin{titlepage}\centering
{\Large\bfseries Universidad de Sonora\par}
\vspace{1.5cm}
{\huge\bfseries\color{marino} EmprendeBúho\par}
\vspace{0.8cm}
{\LARGE\bfseries #1\par}
\vspace{0.5cm}
{\large Ingeniería de Software II\par}
\vspace{0.5cm}
{\large Equipo \guillemotleft Descompuestos por Cholesky\guillemotright\par}
\vspace{0.6cm}
\textbf{Profesor:}\par Adrian Vazquez Osorio\par
\vspace{0.5cm}
\textbf{Integrantes:}\par
Paulina Bojorquez Trujillo\par
Daniel Leinad Domínguez Calvario\par
María Fernanda Hernández García\par
Marco Antonio Tadeo Munro Flores\par
Mario Alberto Ocejo Quijada\par
Angel David Ortega Felix\par
Orlando Pacheco Velasquez\par
Denzel Omar Rivera Urías\par
\vfill
#2\par
Hermosillo, Sonora, 13 de noviembre de 2026\par
\vspace{0.3cm}
{\small Borrador documental para validación del equipo\par}
\end{titlepage}\hypersetup{pageanchor=true}}
\newcommand{\inicioDia}[1]{%
\clearpage\subsection*{Fecha: #1 de 2026}
{\small En \guillemotleft ayer\guillemotright{} se registra la jornada hábil anterior; los lunes se retoma
el viernes, sin atribuir trabajo al fin de semana.}\par
\begingroup\footnotesize
\begin{longtable}{@{}L{2.05cm}L{4.55cm}L{5.05cm}L{4.2cm}@{}}
\rowcolor{dorado}\textcolor{marino}{\textbf{Integrante}}&
\textcolor{marino}{\textbf{¿Qué hiciste ayer?}}&
\textcolor{marino}{\textbf{¿Qué harás hoy?}}&
\textcolor{marino}{\textbf{¿Qué necesitas?}}\\\midrule\endfirsthead
\rowcolor{suave}\textbf{Integrante}&\textbf{¿Qué hiciste ayer?}&
\textbf{¿Qué harás hoy?}&\textbf{¿Qué necesitas?}\\\midrule\endhead
\bottomrule\endfoot}
\newcommand{\finDia}{\end{longtable}\endgroup}
